\chapter*{Tóm tắt}
\addcontentsline{toc}{chapter}{Tóm tắt}

Truy vấn thông tin đa phương thức (Multimodal Information Retrieval) giữa văn bản và hình ảnh truyền thống dựa chủ yếu trên việc so khớp vectơ đặc trưng trong không gian tiềm ẩn chung thông qua các mô hình bộ mã hóa kép như CLIP. Mặc dù đạt hiệu năng cao trong việc nắm bắt tương đồng thị giác bề mặt, cách tiếp cận thuần tương đồng vectơ này bộc lộ hai hạn chế cốt lõi: thứ nhất, mô hình hoàn toàn bỏ qua các tri thức quan hệ có cấu trúc giữa thực thể, thuộc tính và bối cảnh sự kiện; thứ hai, mô hình hoạt động như một hộp đen và không cung cấp cơ chế nội tại giải thích tại sao một kết quả cụ thể được ưu tiên xếp hạng. Đồ án này nghiên cứu và phát triển hệ thống truy vấn đa phương thức tăng cường đồ thị có khả năng giải thích dựa trên cấu trúc quan hệ (Multimodal Retrieval with Graph-based Explainability), được thiết kế theo tiến trình học thuật ba cấp độ (Level 1 $\rightarrow$ Level 2 $\rightarrow$ Level 3) nhằm trả lời câu hỏi về vai trò thực chất của dữ liệu đồ thị và tính trung thực của các giải thích được sinh ra.

Ở Cấp độ 1 (Level 1 -- Graph-enhanced Retrieval), nghiên cứu thiết lập đường cơ sở CLIP đóng băng, xây dựng đồ thị tĩnh từ chú giải đối tượng/scene graph, áp dụng Mạng nơ-ron chú ý đồ thị (GAT) để tính điểm đồ thị và kết hợp rerank với điểm cosine CLIP. Kiểm chứng thực nghiệm trên tập Flickr30k Entities với 3 seed độc lập ghi nhận mô hình kết hợp CLIP + GAT đạt Mean Test MRR $\sim 0.8273$ vượt trội so với CLIP thuần ($0.7232$). Tuy nhiên, phân tích triệt tiêu (ablation study) cho thấy cấu hình không cạnh (adapter với self-loops) cũng đóng góp đáng kể, chứng minh không thể quy toàn bộ phần cải thiện cho cấu trúc quan hệ nếu thiếu kiểm soát thực nghiệm chặt chẽ.

Ở Cấp độ 2 (Level 2 -- Advanced Multimodal Graph-based Retrieval), nghiên cứu mở rộng sang đồ thị tri thức đa phương thức (MMKG) phức tạp trên DB15K và FB15k-237, triển khai cơ chế dung hợp đặc trưng dựa trên AdaMF (AdaMF-based Multimodal Fusion) và phân tích chuyên sâu hiệu năng trên 20 truy vấn đánh giá phân tách rõ nhánh trực tiếp và gián tiếp. Bằng chứng audit độc lập trên DB15K (seed 2027) xác nhận mô hình Full Graph đạt Test MRR $0.4145$ (Recall@10 đạt $0.5556$, indirect MRR đạt $0.3164$), cải thiện vượt bậc so với baseline CLIP thuần ($0.2577$) và vượt cấu hình Shuffled Edges v2 ($0.3156$). Báo cáo ghi nhận minh bạch đặc điểm kỹ thuật của giao thức tạo đồ thị hoán đổi cạnh bảo toàn bậc và kiểm soát rò rỉ tập kiểm thử.

Ở Cấp độ 3 (Level 3 -- Explainability \& Transferability), nghiên cứu xây dựng cơ chế trích xuất đường đi và đồ thị con giải thích bắt nguồn trực tiếp từ cơ chế tính điểm xếp hạng, thiết lập quy trình đo lường định lượng tính trung thực (fidelity), tính hợp lệ (validity) và độ thưa (sparsity) trên các trường hợp truy vấn thành công và thất bại điển hình, đồng thời đánh giá khả năng chuyển giao trên tập dữ liệu MMKG tự xây. Báo cáo giữ nguyên tính trung thực học thuật bằng việc phân định ranh giới rõ ràng giữa các kết quả đã được xác minh bằng artifact thực tế và các hạng mục đang trong quá trình hoàn thiện thực nghiệm.

\textbf{Từ khóa:} Truy vấn đa phương thức; Đồ thị tri thức đa phương thức (MMKG); CLIP; Graph Attention Networks (GAT); Dung hợp AdaMF; Khả năng giải thích đồ thị (xAI); Tính trung thực giải thích (Fidelity).
